\section{Sources}

Two sources carry the whole plan, and neither is reproduced here.

\begin{sloppypar}
\paragraph{The seminar.} MIT 18.A34 Mathematical Problem Solving (Putnam Seminar), Fall 2018, taught by Yufei Zhao, problem sets 1 to 11, at \url{https://ocw.mit.edu/courses/18-a34-mathematical-problem-solving-putnam-seminar-fall-2018/pages/assignments/}, published under CC BY-NC-SA 4.0. Four further sets come from the Fall 2007 edition of the same seminar, 18.S34, by Richard Stanley: pigeonhole, greatest integer, roots of polynomials, and limits. This document refers to problems by set and number only, as in \seminar{PS 9 \#14} or \seminar{Roots \#6}, and paraphrases the sets' own prefaces where they say something about difficulty. The statements are in the sets, not here.

\paragraph{The archive.} Kiran Kedlaya's Putnam archive at \url{https://kskedlaya.org/putnam-archive/}, with every paper from 1985 and solutions from 1995. A year's paper is \texttt{2024.pdf} and its solutions \texttt{2024s.pdf}. The solutions are compiled by Manjul Bhargava, Kiran Kedlaya and Lenhard Ng, who ask that they be linked to rather than copied; this document links, and quotes nothing. Problems are named as \putnam{2021 A1}.
\end{sloppypar}

\paragraph{The format.} The exam's shape, timetable and room rules are from the MAA's \emph{Policies, Rules, and Guidelines} for the 87th William Lowell Putnam Mathematical Competition, 5 December 2026, in the version updated 2026-08-31, which Adit supplied as a PDF on 2026-11-09. Its rules are restated here, not quoted. The policy says the MAA may change it and that its own website governs, so the supervisor's instructions on the day come first.

\paragraph{The tags.} The topic tags on the 120 Putnam problems of 2016 to 2025, the tag counts in Section~1, and the ladders in \secref{sec:ladders}{the four ladders} come from a model's reading of the problem statements and of the fifteen seminar sets. They were assigned by statement, not by solving the problems, so a tag is a guess at the technique and can be wrong. When a call is settled blind because the solution used a technique the tag did not suggest, the tag was wrong too, and that is worth a mark on the section.

\paragraph{The book.} \emph{Putnam and Beyond} by R\u{a}zvan Gelca and Titu Andreescu was used in one way only: to check that the technique sections above cover its parts (methods of proof, algebra, real analysis, geometry and trigonometry, number theory, combinatorics and probability). No text, problem or solution from it is reproduced. A copy of the book that circulates from a WordPress site is not the publisher's.

\paragraph{The examples.} Every worked example in this document was invented for it. None is a competition or textbook problem with its numbers changed. The standard theorems and their standard proofs are the common property of the subject.

